% SLICECANCELGUARD1 2026-10-17 20:20 Asia/Seoul
\documentclass[11pt]{article}
\usepackage{amsmath,amssymb,amsthm,booktabs,hyperref}
\usepackage[margin=.9in]{geometry}
\newtheorem{theorem}{Theorem}
\newtheorem{lemma}[theorem]{Lemma}
\newtheorem{corollary}[theorem]{Corollary}
\title{SLICECANCELGUARD1: Cancellation-Order and Polynomial-Degree Guards for the Exact UNSAT Slice}
\author{P versus NP project}
\date{17 October 2026}
\begin{document}\maketitle
\noindent\textbf{Status.} The cancellation-order counterexample, shell-degree theorem, and moment-duality lemma are \textbf{PROVED/EXACT} at their stated models, pending independent review. The fixed-weight monotone-KW import remains \textbf{OPEN}. No unrestricted set-circuit lower bound or $P\ne NP$ proof is obtained.

\section{IDEAS PROPOSED BEFORE LITERATURE}
A: use the growing order of the exact inclusion--exclusion/factorial-moment series as a set-circuit lower-bound invariant. Falsifier: a small set circuit whose canonical union series already needs growing-order cancellation. B: use exact or harmonic polynomial degree on the slice. Falsifier: the degree lower bound is caused by the threshold shell or yields only linear/trivial circuit bounds. C: import a known hard fixed-weight monotone relation, especially clique/coloring or perfect matching, directly into the full KW relation of the exact UNSAT slice. Missing theorem: a model-matched fixed-weight monotone lower bound plus an output-preserving encoding.

\section{LITERATURE CHECK AND REVISIONS}
Kahn--Linial--Samorodnitsky study how many intersection moments are needed to approximate a union and show that low-order inclusion--exclusion information can be insufficient in general. Cluster/cumulant methods similarly encode non-existence probabilities through connected dependencies. These support A as a diagnostic, not a gate-stable circuit invariant.

Filmus--Ihringer and later Filmus/Poupin prove junta theorems for low-degree Boolean functions on a Hamming slice; Chiarelli--Hatami--Saks show a real degree-$d$ Boolean function on the cube depends on at most $C2^d$ relevant coordinates. These results make polynomial degree a legitimate structural control, but do not give superpolynomial monotone circuit lower bounds.

A targeted search found no theorem giving an unrestricted exponential monotone Boolean-circuit lower bound for perfect matching specifically on fixed-edge-count or regular/cubic input slices. AC0-equivalence of general and cubic matching is not a monotone-circuit-preserving lower-bound transfer. Known exponential results for multilinear/consistent circuits and exact-matching extension complexity are restricted-model results and are not imported.

\section{CONSTRUCTION OR PROOF ATTEMPT}
Let $\Omega=\binom{[N]}t$ and $Y_i=\{X\in\Omega:i\in X\}$.

\begin{theorem}[Cancellation order is not a set-circuit size measure]
For any $1\le m\le N$, the family
\[S_m=\bigcup_{i=1}^m Y_i\]
has a binary union set circuit of size $m-1$. Under the uniform exact-$t$ measure, its canonical inclusion--exclusion expansion is
\[\Pr[S_m]=\sum_{j=1}^{\min\{m,t\}}(-1)^{j+1}\binom mj\frac{(t)_j}{(N)_j},\]
and every displayed term is nonzero. Thus a circuit of linear size can have nonzero cancellation terms through order $\min\{m,t\}$.
\end{theorem}
\begin{proof}
The circuit is a balanced or sequential OR of the $m$ coordinate stars. Every $j$ distinct stars intersect in the family of exact-$t$ sets containing those $j$ coordinates, whose measure is $(t)_j/(N)_j$ for $j\le t$ and zero afterward.
\end{proof}

\begin{corollary}
The proposed implication ``small union/intersection circuit implies bounded inclusion--exclusion or factorial-cumulant order'' is \textbf{FALSIFIED/RETIRED}. In particular, the growing-order cancellation observed for the random-SAT witness count cannot by itself imply a superpolynomial set-circuit lower bound.
\end{corollary}

Now let $F:\{0,1\}^N\to\{0,1\}$ be any total slice extension with threshold $t$: $F=0$ below weight $t$, $F=1$ above weight $t$, and arbitrary on weight $t$. Let
\[\alpha=\binom Nt^{-1}\sum_{|x|=t}F(x).\]

\begin{theorem}[Universal threshold-shell degree]
If $0<\alpha<1$, every real multilinear polynomial representing $F$ on the Boolean cube has degree at least $N-1$.
\end{theorem}
\begin{proof}
Average a representing polynomial over all coordinate permutations. The resulting symmetric polynomial has no larger degree and, on Boolean inputs of weight $w$, equals a univariate polynomial $p(w)$ with
\[p(0)=\cdots=p(t-1)=0,\quad p(t)=\alpha,\quad p(t+1)=\cdots=p(N)=1.\]
The finite difference $\Delta p(w)=p(w+1)-p(w)$ is nonzero at $w=t-1,t$ and zero at each of the other $N-2$ integer points $w\in\{0,\ldots,N-1\}$. Since $\Delta p$ is a nonzero polynomial of degree at most $\deg p-1$, it must have degree at least $N-2$, hence $\deg p\ge N-1$.
\end{proof}

For the exact-density UNSAT slice, both SAT and UNSAT formulas exist on the critical layer for all sufficiently large parameters, so $0<\alpha<1$. Hence its total Valiant extension has degree at least $N-1$.

\begin{corollary}[Shell-degree hostile control]
The degree bound above is not SAT-specific: it holds for every nontrivial slice, including computationally easy slices. Therefore almost-maximal full-cube degree cannot be the load-bearing lower-bound invariant for the internal exact-layer function.
\end{corollary}

\begin{lemma}[Moment matching is a low-degree dual certificate]
Let $\mu_+$ and $\mu_-$ be distributions on the exact layer. If
\[\mathbb E_{\mu_+}\prod_{i\in A}x_i=\mathbb E_{\mu_-}\prod_{i\in A}x_i\]
for every $|A|\le r$, then every real polynomial of degree at most $r$ has the same expectation under $\mu_+$ and $\mu_-$. Thus an exact moment-matching UNSAT/SAT pair is a dual certificate against all degree-$r$ polynomial distinguishers.
\end{lemma}
\begin{proof}
Multilinearize the polynomial and expand it in monomials $\prod_{i\in A}x_i$.
\end{proof}

This recovers the HD/orbit-moment program as a genuine low-degree lower-bound mechanism, but low polynomial degree is not known to characterize polynomial-size monotone set circuits. It remains an approximation layer rather than the final size invariant.

\section{FIXED-WEIGHT KW IMPORT AUDIT}
The most promising attempted escape is to reduce a known hard monotone KW relation into the full exact-slice KW relation without mapping the whole Boolean cube into one constant-weight layer. A relation-level reduction only needs positive and negative source families plus an output map.

A targeted literature search found unrestricted monotone hardness for full matching/clique functions and restricted-model exponential results, but no directly usable theorem asserting exponential unrestricted monotone Boolean-circuit hardness on a fixed-edge-count or regular/cubic slice. Perfect matching on cubic graphs is AC0-equivalent to general matching as a decision/search problem, but that equivalence is not known here to preserve the monotone Boolean circuit model or the directional KW witness relation. Consequently this route remains \textbf{OPEN}, not adopted.

\section{TESTS AND OUTCOME}
The main positive result is conceptual narrowing. The infinite factorial-moment series remains useful for understanding random-SAT dependence and for constructing moment-matching distributions, but its cancellation depth cannot serve as a gate count. Full-cube polynomial degree is nearly maximal for a trivial reason shared by all nontrivial total slices. The only surviving uses of the series/HD machinery are (i) bounded-width DNF fooling via Bonferroni, and (ii) low-degree dual certificates inside a larger approximation theorem.

The highest-value next targets are now: (1) a gate-stable set-circuit invariant not reducible to cancellation order, raw degree, minterm count, or additive HD mass; and (2) a genuine fixed-weight monotone KW lower bound that can be embedded into the exact UNSAT-vs-SAT relation. The full KW relation remains preferable to assignment-masked restrictions, which already have polynomial DAG protocols.

Root main.tex remains truncated and was not modified.

\begin{thebibliography}{7}
\bibitem{KLS96} J. Kahn, N. Linial, A. Samorodnitsky, Inclusion-exclusion: Exact and approximate, Combinatorica 16, 1996.
\bibitem{Filmus18} Y. Filmus, F. Ihringer, Boolean constant degree functions on the slice are juntas, 2018.
\bibitem{Filmus22} Y. Filmus, Junta threshold for low degree Boolean functions on the slice, 2022.
\bibitem{CHS20} J. Chiarelli, P. Hatami, M. Saks, An Asymptotically Tight Bound on the Number of Relevant Variables in a Bounded Degree Boolean Function, Combinatorica 40, 2020.
\bibitem{PV04} A. K. Ponnuswami, H. Venkateswaran, Monotone multilinear Boolean circuits for bipartite perfect matching require exponential size, FSTTCS 2004.
\bibitem{JiaSY23} X. Jia, O. Svensson, W. Yuan, The Exact Bipartite Matching Polytope Has Exponential Extension Complexity, SODA 2023.
\bibitem{Valiant86} L. G. Valiant, Negation Is Powerless for Boolean Slice Functions, SIAM J. Comput. 15(2), 1986.
\end{thebibliography}
\end{document}
% FINAL_VERIFICATION_SENTINEL::MAIN::SLICECANCELGUARD1::2026-10-17_2020::P_VS_NP_OPEN